\subsection{\texorpdfstring{ECCLESIOLOGICAL AND CANONICAL
CONSEQUENCES~\\
OF THE SACRAMENTAL NATURE OF THE CHURCH:\\
ECCLESIAL COMMUNION, CONCILIARITY AND
AUTHORITY}{ECCLESIOLOGICAL AND CANONICAL CONSEQUENCES~ OF THE SACRAMENTAL NATURE OF THE CHURCH: ECCLESIAL COMMUNION, CONCILIARITY AND AUTHORITY}}\label{ecclesiological-and-canonical-consequences-of-the-sacramental-nature-of-the-church-ecclesial-communion-conciliarity-and-authority}

\subsubsection{Ravenna, 13 October 2007}\label{ravenna-13-october-2007}

\subsubsection{~}\label{section}

\subsection{\texorpdfstring{\textbf{\textbf{{~}}}}{~}}\label{section-1}

\subsection{Introduction}\label{introduction}

1. ``That they may all be one. As you, Father, are in me and I am in
you, may they also be one in us so that the world may believe that you
have sent me'' (Jn 17, 21). We give thanks to the triune God who has
gathered us -- members of the Joint International Commission for the
Theological Dialogue between the Roman Catholic Church and the Orthodox
Church - so that we might respond together in obedience to this prayer
of Jesus. We are conscious that our dialogue is restarting in a world
that has changed profoundly in recent times. The processes of
secularization and globalization, and the challenge posed by new
encounters between Christians and believers of other religions, require
that the disciples of Christ give witness to their faith, love and hope
with a new urgency. May the Spirit of the risen Lord empower our hearts
and minds to bear the fruits of unity in the relationship between our
Churches, so that together we may serve the unity and peace of the whole
human family. May the same Spirit lead us to the full expression of the
mystery of ecclesial communion, that we gratefully acknowledge as a
wonderful gift of God to the world, a mystery whose beauty radiates
especially in the holiness of the saints, to which all are called.

2. Following the plan adopted at its first meeting in Rhodes in 1980,
the Joint Commission began by addressing the mystery of
ecclesial~\emph{koinônia}~in the light of the mystery of the Holy
Trinity and of the Eucharist. This enabled a deeper understanding of
ecclesial communion, both at the level of the local community around its
bishop, and at the level of relations between bishops and between the
local Churches over which each presides in communion with the One Church
of God extending across the universe (cfr. Munich Document, 1982). In
order to clarify the nature of communion, the Joint Commission
underlined the relationship which exists between faith, the sacraments
-- especially the three sacraments of Christian initiation -- and the
unity of the Church (cfr. Bari\emph{~}Document, 1987). Then by studying
the sacrament of Order in the sacramental structure of the Church, the
Commission\emph{~}indicated clearly the role of apostolic succession as
the guarantee of the~\emph{koinonia}~of the whole Church and of its
continuity with the Apostles in every time and place (cfr. Valamo
Document, 1988). From 1990 until 2000, the main subject discussed by the
Commission was that of ``uniatism'' (Balamand Document, 1993; Baltimore,
2000), a subject to which we shall give further consideration in the
near future. Now we take up the theme raised at the end of the Valamo
Document, and reflect upon ecclesial communion, conciliarity and
authority.

3. On the basis of these common affirmations of our faith, we must now
draw the ecclesiological and canonical consequences which flow from the
sacramental nature of the Church. Since the Eucharist, in the light of
the Trinitarian mystery, constitutes the criterion of ecclesial life as
a whole, how do institutional structures visibly reflect the mystery of
this~\emph{koinonia}? Since the one and holy Church is realised both in
each local Church celebrating the Eucharist and at the same time in
the~\emph{koinonia}~of all the Churches, how does the life of the
Churches manifest this sacramental structure?

4. Unity and multiplicity, the relationship between the one Church and
the many local Churches, that constitutive relationship of the Church,
also poses the question of the relationship between the authority
inherent in every ecclesial institution and the conciliarity which flows
from the mystery of the Church as communion. As the terms ``authority''
and ``conciliarity'' cover a very wide area, we shall begin by defining
the way we understand them.\phantomsection\label{_ftnref1}{}{[}1{]}

~

\subsection{I. The Foundations of Conciliarity and of
Authority}\label{i.-the-foundations-of-conciliarity-and-of-authority}

\subsubsection{1. Conciliarity}\label{conciliarity}

5. The term conciliarity or synodality comes from the word ``council''
(\emph{synodos}~in Greek,~\emph{concilium}~in Latin), which primarily
denotes a gathering of bishops exercising a particular responsibility.
It is also possible, however, to take the term in a more comprehensive
sense referring to all the members of the Church (cfr. the Russian
term~\emph{sobornost}). Accordingly we shall speak first of all of
conciliarity as signifying that each member of the Body of Christ, by
virtue of baptism, has his or her place and proper responsibility in
eucharistic~\emph{koinonia}(\emph{communio}~in Latin). Conciliarity
reflects the Trinitarian mystery and finds therein its ultimate
foundation. The three persons of the Holy Trinity are ``enumerated'', as
St Basil the Great says (\emph{On the Holy Spirit}, 45), without the
designation as ``second'' or ``third'' person implying any diminution or
subordination. Similarly, there also exists an order (\emph{taxis})
among local Churches, which however does not imply inequality in their
ecclesial nature.

6. The Eucharist manifests the Trinitarian~\emph{koinônia}~actualized in
the faithful as an organic unity of several members each of whom has a
charism, a service or a proper ministry, necessary in their variety and
diversity for the edification of all in the one ecclesial Body of Christ
(cfr. 1 Cor 12, 4-30). All are called, engaged and held accountable --
each in a different though no less real manner -- in the common
accomplishment of the actions which, through the Holy Spirit, make
present in the Church the ministry of Christ, ``the way, the truth and
the life'' (Jn 14, 6). In this way, the mystery of
salvific~\emph{koinonia}~with the Blessed Trinity is realized in
humankind.

7. The whole community and each person in it bears the ``conscience of
the Church'' (\emph{ekklesiastike syneidesis}), as Greek theology calls
it, the~\emph{sensus fidelium}~in Latin terminology. By virtue of
Baptism and Confirmation (Chrismation) each member of the Church
exercises a form of authority in the Body of Christ. In this sense, all
the faithful (and not just the bishops) are responsible for the faith
professed at their Baptism. It is our common teaching that the people of
God, having received ``the anointing which comes from the Holy One'' (1
Jn 2, 20 and 27), in communion with their pastors, cannot err in matters
of faith (cfr. Jn 16, 13).

8. In proclaiming the Church's faith and in clarifying the norms of
Christian conduct, the bishops have a specific task by divine
institution. ``As successors of the Apostles, the bishops are
responsible for communion in the apostolic faith and for fidelity to the
demands of a life in keeping with the Gospel'' (Valamo Document, n. 40).

9. Councils are the principal way in which communion among bishops is
exercised (cfr. Valamo Document, n. 52). For ``attachment to the
apostolic communion binds all the bishops together linking
the~\emph{episkope~}of the local Churches to the College of the
Apostles. They too form a college rooted by the Spirit in the `once for
all' of the apostolic group, the unique witness to the faith. This means
not only that they should be united among themselves in faith, charity,
mission, reconciliation, but that they have in common the same
responsibility and the same service to the Church'' (Munich Document,
III, 4).

10. This conciliar dimension of the Church's life belongs to its
deep-seated nature. That is to say, it is founded in the will of Christ
for his people (cfr. Mt 18, 15-20), even if its canonical realizations
are of necessity also determined by history and by the social, political
and cultural context. Defined thus, the conciliar dimension of the
Church is to be found at the three levels of ecclesial communion, the
local, the regional and the universal: at the local level of the diocese
entrusted to the bishop; at the regional level of a group of local
Churches with their bishops who ``recognize who is the first amongst
themselves'' (Apostolic Canon 34); and at the universal level, where
those who are first (\emph{protoi}) in the various regions, together
with all the bishops, cooperate in that which concerns the totality of
the Church. At this level also, the~\emph{protoi}~must recognize who is
the first amongst themselves.

11. The Church exists in many and different places, which manifests its
catholicity. Being ``catholic'', it is a living organism, the Body of
Christ. Each local Church, when in communion with the other local
Churches, is a manifestation of the one and indivisible Church of God.
To be ``catholic'' therefore means to be in communion with the one
Church of all times and of all places. That is why the breaking of
eucharistic communion means the wounding of one of the essential
characteristics of the Church, its catholicity.

\subsubsection{2. Authority}\label{authority}

12. When we speak of authority, we are referring to~\emph{exousia}, as
it is described in the New Testament. The authority of the Church comes
from its Lord and Head, Jesus Christ. Having received his authority from
God the Father, Christ after his Resurrection shared it, through the
Holy Spirit, with the Apostles (cfr. Jn 20, 22). Through the Apostles it
was transmitted to the bishops, their successors, and through them to
the whole Church. Jesus Christ our Lord exercised this authority in
various ways whereby, until its eschatological fulfilment (cfr. 1 Cor
15, 24-28), the Kingdom of God manifests itself to the world: by
teaching (cfr. Mt 5, 2; Lk 5, 3); by performing miracles (cfr. Mk 1,
30-34; Mt 14, 35-36); by driving out impure spirits (cfr. Mk 1, 27; Lk
4, 35-36); in the forgiveness of sins (cfr. Mk 2, 10; Lk 5, 24); and in
leading his disciples in the ways of salvation (cfr. Mt 16, 24). In
conformity with the mandate received from Christ (cfr. Mt 28, 18-20),
the exercise of the authority proper to the apostles and afterwards to
the bishops includes the proclamation and the teaching of the Gospel,
sanctification through the sacraments, particularly the Eucharist, and
the pastoral direction of those who believe (cfr. Lk 10, 16).

13. Authority in the Church belongs to Jesus Christ himself, the one
Head of the Church (cfr. Eph 1, 22; 5, 23). By his Holy Spirit, the
Church as his Body shares in his authority (cfr. Jn 20, 22-23).
Authority in the Church has as its goal the gathering of the whole of
humankind into Jesus Christ (cfr. Eph 1,10; Jn 11, 52). The authority
linked with the grace received in ordination is not the private
possession of those who receive it nor something delegated from the
community; rather, it is a gift of the Holy Spirit destined for the
service (\emph{diakonia}) of the community and never exercised outside
of it. Its exercise includes the participation of the whole community,
the bishop being in the Church and the Church in the bishop (cfr. St
Cyprian,~\emph{Ep}. 66, 8).

14. The exercise of authority accomplished in the Church, in the name of
Christ and by the power of the Holy Spirit, must be, in all its forms
and at all levels, a service (\emph{diakonia}) of love, as was that of
Christ (cfr. Mk 10, 45; Jn 13, 1-16). The authority of which we are
speaking, since it expresses divine authority, cannot subsist in the
Church except in the love between the one who exercises it and those
subject to it. It is, therefore, an authority without domination,
without physical or moral coercion. Since it is a participation in
the~\emph{exousia}~of the crucified and exalted Lord, to whom has been
given all authority in heaven and on earth (cfr. Mt 28, 18), it can and
must call for obedience. At the same time, because of the Incarnation
and the Cross, it is radically different from that of leaders of nations
and of the great of this world (cfr. Lk 22, 25-27). While this authority
is certainly entrusted to people who, because of weakness and sin, are
often tempted to abuse it, nevertheless by its very nature the
evangelical identification between authority and service constitutes a
fundamental norm for the Church. For Christians, to rule is to serve.
The exercise and spiritual efficacy of ecclesial authority are thereby
assured through free consent and voluntary co-operation. At a personal
level, this translates into obedience to the authority of the Church in
order to follow Christ who was lovingly obedient to the Father even unto
death and death on a Cross (cfr. Phil 2, 8).

15. Authority within the Church is founded upon the Word of God, present
and alive in the community of the disciples. Scripture is the revealed
Word of God, as the Church, through the Holy Spirit present and active
within it, has discerned it in the living Tradition received from the
Apostles. At the heart of this Tradition is the Eucharist (cfr. 1 Cor
10, 16-17; 11, 23-26). The authority of Scripture derives from the fact
that it is the Word of God which, read in the Church and by the Church,
transmits the Gospel of salvation. Through Scripture, Christ addresses
the assembled community and the heart of each believer. The Church,
through the Holy Spirit present within it, authentically interprets
Scripture, responding to the needs of times and places. The constant
custom of the Councils to enthrone the Gospels in the midst of the
assembly both attests the presence of Christ in his Word, which is the
necessary point of reference for all their discussions and decisions,
and at the same time affirms the authority of the Church to interpret
this Word of God.

16. In his divine Economy, God wills that his Church should have a
structure oriented towards salvation. To this essential structure belong
the faith professed and the sacraments celebrated in the apostolic
succession. Authority in the ecclesial communion is linked to this
essential structure: its exercise is regulated by the canons and
statutes of the Church. Some of these regulations may be differently
applied according to the needs of ecclesial communion in different times
and places, provided that the essential structure of the Church is
always respected. Thus, just as communion in the sacraments presupposes
communion in the same faith (cfr. Bari Document, nn.29-33), so too, in
order for there to be full ecclesial communion, there must be, between
our Churches, reciprocal recognition of canonical legislations in their
legitimate diversities.

\subsection{\texorpdfstring{\textbf{~}}{~}}\label{section-2}

\subsection{II. The threefold actualization of Conciliarity and
Authority}\label{ii.-the-threefold-actualization-of-conciliarity-and-authority}

17. Having pointed out the foundation of conciliarity and of authority
in the Church, and having noted the complexity of the content of these
terms, we must now reply to the following questions: How do
institutional elements of the Church visibly express and serve the
mystery of~\emph{koinonia}? How do the canonical structures of the
Churches express their sacramental life? To this end we distinguished
between three levels of ecclesial institutions: that of the local Church
around its bishop; that of a region taking in several neighbouring local
Churches; and that of the whole inhabited earth (\emph{oikoumene}) which
embraces all the local Churches.

\subsubsection{1. The Local Level}\label{the-local-level}

18. The Church of God exists where there is a community gathered
together in the Eucharist, presided over, directly or through his
presbyters, by a bishop legitimately ordained into the apostolic
succession, teaching the faith received from the Apostles, in communion
with the other bishops and their Churches. The fruit of this Eucharist
and this ministry is to gather into an authentic communion of faith,
prayer, mission, fraternal love and mutual aid, all those who have
received the Spirit of Christ in Baptism. This communion is the frame in
which all ecclesial authority is exercised. Communion is the criterion
for its exercise.

19. Each local Church has as its mission to be, by the grace of God, a
place where God is served and honoured, where the Gospel is announced,
where the sacraments are celebrated, where the faithful strive to
alleviate the world's misery, and where each believer can find
salvation. It is the light of the world (cfr. Mt 5, 14-16), the leaven
(cfr. Mt 13, 33), the priestly community of God (cfr. 1 Pet 2, 5 and 9).
The canonical norms which govern it aim at ensuring this mission.

20. By virtue of that very Baptism which made him or her a member of
Christ, each baptized person is called, according to the gifts of the
one Holy Spirit, to serve within the community (cfr. 1 Cor 12, 4-27).
Thus through communion, whereby all the members are at the service of
each other, the local Church appears already ``synodal'' or
``conciliar'' in its structure. This ``synodality'' does not show itself
only in the relationships of solidarity, mutual assistance and
complementarity which the various ordained ministries have among
themselves. Certainly, the presbyterium is the council of the bishop
(cfr. St Ignatius of Antioch,~\emph{To the Trallians}, 3), and the
deacon is his ``right arm'' (\emph{Didascalia Apostolorum}, 2, 28, 6),
so that, according to the recommendation of St Ignatius of Antioch,
everything be done in concert (cfr.~\emph{To the Ephesians}, 6).
Synodality, however, also involves all the members of the community in
obedience to the bishop, who is the~\emph{protos}~and head
(\emph{kephale}) of the local Church, required by ecclesial communion.
In keeping with Eastern and Western traditions, the active participation
of the laity, both men and women, of monastics and consecrated persons,
is effected in the diocese and the parish through many forms of service
and mission.

21. The charisms of the members of the community have their origin in
the one Holy Spirit, and are directed to the good of all. This fact
sheds light on both the demands and the limits of the authority of each
one in the Church. There should be neither passivity nor substitution of
functions, neither negligence nor domination of anyone by another. All
charisms and ministries in the Church converge in unity under the
ministry of the bishop, who serves the communion of the local Church.
All are called to be renewed by the Holy Spirit in the sacraments and to
respond in constant repentance (\emph{metanoia}), so that their
communion in truth and charity is ensured.

\subsubsection{2. The Regional Level}\label{the-regional-level}

22. Since the Church reveals itself to be catholic in
the~\emph{synaxis}~of the local Church, this catholicity must truly
manifest itself in communion with the other Churches which confess the
same apostolic faith and share the same basic ecclesial structure,
beginning with those close at hand in virtue of their common
responsibility for mission in that region which is theirs (cfr. Munich
Document, III, 3, and Valamo Document, nn.52 and 53). Communion among
Churches is expressed in the ordination of bishops. This ordination is
conferred according to canonical order by three or more bishops, or at
least two (cfr. Nicaea I, Canon 4), who act in the name of the episcopal
body and of the people of God, having themselves received their ministry
from the Holy Spirit by the imposition of hands in the apostolic
succession. When this is accomplished in conformity with the canons,
communion among Churches in the true faith, sacraments and ecclesial
life is ensured, as well as living communion with previous generations.

23. Such effective communion among several local Churches, each being
the Catholic Church in a particular place, has been expressed by certain
practices: the participation of the bishops of neighbouring sees at the
ordination of a bishop to the local Church; the invitation to a bishop
from another Church to concelebrate at the~\emph{synaxis}~of the local
Church; the welcome extended to the faithful from these other Churches
to partake of the eucharistic table; the exchange of letters on the
occasion of an ordination; and the provision of material assistance.

24. A canon accepted in the East as in the West, expresses the
relationship between the local Churches of a region: ``The bishops of
each province (\emph{ethnos}) must recognize the one who is first
(\emph{protos}) amongst them, and consider him to be their head
(\emph{kephale}), and not do anything important without his consent
(\emph{gnome}); each bishop may only do what concerns his own diocese
(\emph{paroikia}) and its dependent territories. But the first
(\emph{protos}) cannot do anything without the consent of all. For in
this way concord (\emph{homonoia}) will prevail, and God will be praised
through the Lord in the Holy Spirit'' (Apostolic Canon 34).

25. This norm, which re-emerges in several forms in canonical tradition,
applies to all the relations between the bishops of a region, whether
those of a province, a metropolitanate, or a patriarchate. Its practical
application may be found in the synods or the councils of a province,
region or patriarchate. The fact that the composition of a regional
synod is always essentially episcopal, even when it includes other
members of the Church, reveals the nature of synodal authority. Only
bishops have a deliberative voice. The authority of a synod is based on
the nature of the episcopal ministry itself, and manifests the collegial
nature of the episcopate at the service of the communion of Churches.

26. A synod (or council) in itself implies the participation of all the
bishops of a region. It is governed by the principle of consensus and
concord (\emph{homonoia}), which is signified by eucharistic
concelebration, as is implied by the final doxology of the
above-mentioned Apostolic Canon 34. The fact remains, however, that each
bishop in his pastoral care is judge, and is responsible before God for
the affairs of his own diocese (cfr. St Cyprian,~\emph{Ep.}55, 21); thus
he is the guardian of the catholicity of his local Church, and must be
always careful to promote catholic communion with other Churches.

27. It follows that a regional synod or council does not have any
authority over other ecclesiastical regions. Nevertheless, the exchange
of information and consultations between the representatives of several
synods are a manifestation of catholicity, as well as of that fraternal
mutual assistance and charity which ought to be the rule between all the
local Churches, for the greater common benefit. Each bishop is
responsible for the whole Church together with all his colleagues in one
and the same apostolic mission.

28. In this manner several ecclesiastical provinces have come to
strengthen their links of common responsibility. This was one of the
factors giving rise to the patriarchates in the history of our Churches.
Patriarchal synods are governed by the same ecclesiological principles
and the same canonical norms as provincial synods.

29. In subsequent centuries, both in the East and in the West, certain
new configurations of communion between local Churches have developed.
New patriarchates and autocephalous Churches have been founded in the
Christian East, and in the Latin Church there has recently emerged a
particular pattern of grouping of bishops, the Episcopal Conferences.
These are not, from an ecclesiological standpoint, merely administrative
subdivisions: they express the spirit of communion in the Church, while
at the same time respecting the diversity of human cultures.

30. In fact, regional synodality, whatever its contours and canonical
regulation, demonstrates that the Church of God is not a communion of
persons or local Churches cut off from their human roots. Because it is
the community of salvation and because this salvation is ``the
restoration of creation'' (cfr. St Irenaeus,~\emph{Adv. Haer}., 1, 36,
1), it embraces the human person in everything which binds him\emph{~}or
her to human reality as created by God. The Church is not just a
collection of individuals; it is made up of communities with different
cultures, histories and social structures.

31. In the grouping of local Churches at the regional level, catholicity
appears in its true light. It is the expression of the presence of
salvation not in an undifferentiated universe but in humankind as God
created it and comes to save it. In the mystery of salvation, human
nature is at the same time both assumed in its fullness and cured of
what sin has infused into it by way of self-sufficiency, pride, distrust
of others, aggressiveness, jealousy, envy, falsehood and hatred.
Ecclesial~\emph{koinonia}~is the gift by which all humankind is joined
together, in the Spirit of the risen Lord. This unity, created by the
Spirit, far from lapsing into uniformity, calls for and thus preserves
-- and, in a certain way, enhances -- diversity and particularity.

\subsubsection{3. The Universal Level}\label{the-universal-level}

32. Each local Church is in communion not only with neighbouring
Churches, but with the totality of the local Churches, with those now
present in the world, those which have been since the beginning, and
those which will be in the future, and with the Church already in glory.
According to the will of Christ, the Church is one and indivisible, the
same always and in every place. Both sides confess, in the
Nicene-Constantinopolitan Creed, that the Church is one and catholic.
Its catholicity embraces not only the diversity of human communities but
also their fundamental unity.

33. It is clear, therefore, that one and the same faith is to be
confessed and lived out in all the local Churches, the same unique
Eucharist is to be celebrated everywhere, and one and the same apostolic
ministry is to be at work in all the communities. A local Church cannot
modify the Creed, formulated by the ecumenical Councils, although the
Church ought always ``to give suitable answers to new problems, answers
based on the Scriptures and in accord and essential continuity with the
previous expressions of dogmas'' (Bari Document, n.29). Equally, a local
Church cannot change a fundamental point regarding the form of ministry
by a unilateral decision, and no local Church can celebrate the
Eucharist in wilful separation from other local Churches without
seriously affecting ecclesial communion. In all of these things one
touches on the bond of communion itself -- thus, on the very being of
the Church.

34. It is because of this communion that all the Churches, through
canons, regulate everything relating to the Eucharist and the
sacraments, the ministry and ordination, and the handing on
(\emph{paradosis}) and teaching (\emph{didaskalia}) of the faith. It is
clear why in this domain canonical rules and disciplinary norms are
needed.

35. In the course of history, when serious problems arose affecting the
universal communion and concord between Churches -- in regard either to
the authentic interpretation of the faith, or to ministries and their
relationship to the whole Church, or to the common discipline which
fidelity to the Gospel requires - recourse was made to Ecumenical
Councils. These councils were ecumenical not just because they assembled
together bishops from all regions and particularly those of the five
major sees, Rome, Constantinople, Alexandria, Antioch and Jerusalem,
according to the ancient order (\emph{taxis}). It was also because their
solemn doctrinal decisions and their common faith formulations,
especially on crucial points, are binding for all the Churches and all
the faithful, for all times and all places. This is why the decisions of
the Ecumenical Councils remain normative.

36. The history of the Ecumenical Councils shows what are to be
considered their special characteristics. This matter needs to be
studied further in our future dialogue, taking account of the evolution
of ecclesial structures during recent centuries in the East and the
West.

37. The ecumenicity of the decisions of a council is recognized through
a process of reception of either long or short duration, according to
which the people of God as a whole -- by means of reflection,
discernment, discussion and prayer - acknowledge in these decisions the
one apostolic faith of the local Churches, which has always been the
same and of which the bishops are the teachers (\emph{didaskaloi}) and
the guardians. This process of reception is differently interpreted in
East and West according to their respective canonical traditions.

38. Conciliarity or synodality involves, therefore, much more than the
assembled bishops. It involves also their Churches. The former are
bearers of and give voice to the faith of the latter. The bishops'
decisions have to be received in the life of the Churches, especially in
their liturgical life. Each Ecumenical Council received as such, in the
full and proper sense, is, accordingly, a manifestation of and service
to the communion of the whole Church.

39. Unlike diocesan and regional synods, an Ecumenical Council is not an
``institution'' whose frequency can be regulated by canons; it is rather
an ``event'', a~\emph{kairos}~inspired by the Holy Spirit who guides the
Church so as to engender within it the institutions which it needs and
which respond to its nature. This harmony between the Church and the
councils is so profound that, even after the break between East and West
which rendered impossible the holding of Ecumenical Councils in the
strict sense of the term, both Churches continued to hold councils
whenever serious crises arose. These councils gathered together the
bishops of local Churches in communion with the See of Rome or, although
understood in a different way, with the See of Constantinople,
respectively. In the Roman Catholic Church, some of these councils held
in the West were regarded as ecumenical. This situation, which obliged
both sides of Christendom to convoke councils proper to each of them,
favoured dissensions which contributed to mutual estrangement. The means
which will allow the re-establishment of ecumenical consensus must be
sought out.

40. During the first millennium, the universal communion of the Churches
in the ordinary course of events was maintained through fraternal
relations between the bishops. These relations, among the bishops
themselves, between the bishops and their respective~\emph{protoi}, and
also among the~\emph{protoi}~themselves in the canonical order
(\emph{taxis}) witnessed by the ancient Church, nourished and
consolidated ecclesial communion. History records the consultations,
letters and appeals to major sees, especially to that of Rome, which
vividly express the solidarity that~\emph{koinonia~}creates. Canonical
provisions such as the inclusion of the names of the bishops of the
principal sees in the diptychs and the communication of the profession
of faith to the other patriarchs on the occasion of elections, are
concrete expressions of~\emph{koinonia}.

41. Both sides agree that this canonical~\emph{taxis}~was recognised by
all in the era of the undivided Church. Further, they agree that Rome,
as the Church that ``presides in love'' according to the phrase of St
Ignatius of Antioch (\emph{To the Romans}, Prologue), occupied the first
place in the~\emph{taxis}, and that the bishop of Rome was therefore
the~\emph{protos}~among the patriarchs. They disagree, however, on the
interpretation of the historical evidence from this era regarding the
prerogatives of the bishop of Rome as~\emph{protos}, a matter that was
already understood in different ways in the first millennium.

42. Conciliarity at the universal level, exercised in the ecumenical
councils, implies an active role of the bishop of Rome,
as~\emph{protos}~of the bishops of the major sees, in the consensus of
the assembled bishops. Although the bishop of Rome did not convene the
ecumenical councils of the early centuries and never personally presided
over them, he nevertheless was closely involved in the process of
decision-making by the councils.

43. Primacy and conciliarity are mutually interdependent. That is why
primacy at the different levels of the life of the Church, local,
regional and universal, must always be considered in the context of
conciliarity, and conciliarity likewise in the context of primacy.

Concerning primacy at the different levels, we wish to affirm the
following points:

1. Primacy at all levels is a practice firmly grounded in the canonical
tradition of the Church.

2. While the fact of primacy at the universal level is accepted by both
East and West, there are differences of understanding with regard to the
manner in which it is to be exercised, and also with regard to its
scriptural and theological foundations.

44. In the history of the East and of the West, at least until the ninth
century, a series of prerogatives was recognised, always in the context
of conciliarity, according to the conditions of the times, for
the~\emph{protos}~or~\emph{kephale}~at each of the established
ecclesiastical levels: locally, for the bishop as~\emph{protos}~of his
diocese with regard to his presbyters and people; regionally, for
the~\emph{protos}~of each metropolis with regard to the bishops of his
province, and for the~\emph{protos}~of each of the five patriarchates,
with regard to the metropolitans of each circumscription; and
universally, for the bishop of Rome as~\emph{protos}~among the
patriarchs. This distinction of levels does not diminish the sacramental
equality of every bishop or the catholicity of each local Church.

~

\subsection{Conclusion}\label{conclusion}

45. It remains for the question of the role of the bishop of Rome in the
communion of all the Churches to be studied in greater depth. What is
the specific function of the bishop of the ``first see'' in an
ecclesiology of~\emph{koinonia~}and in view of what we have said on
conciliarity and authority in the present text? How should the teaching
of the first and second Vatican councils on the universal primacy be
understood and lived in the light of the ecclesial practice of the first
millennium? These are crucial questions for our dialogue and for our
hopes of restoring full communion between us.

46. We, the members of the Joint International Commission for the
Theological Dialogue between the Roman Catholic Church and the Orthodox
Church, are convinced that the above statement on ecclesial communion,
conciliarity and authority represents positive and significant progress
in our dialogue, and that it provides a firm basis for future discussion
of the question of primacy at the universal level in the Church. We are
conscious that many difficult questions remain to be clarified, but we
hope that, sustained by the prayer of Jesus ``That they may all be one
\ldots{} so that the world may believe'' (Jn 17, 21), and in obedience
to the Holy Spirit, we can build upon the agreement already reached.
Reaffirming and confessing ``one Lord, one faith, one baptism'' (Eph 4,
5), we give glory to God the Holy Trinity, Father, Son and Holy Spirit,
who has gathered us together.

\phantomsection\label{_ftn1}{}{[}1{]}~Orthodox participants felt it
important to emphasize that the use of the terms ``the Church'', ``the
universal Church'', ``the indivisible Church'' and ``the Body of
Christ'' in this document and in similar documents produced by the Joint
Commission in no way undermines the self-understanding of the Orthodox
Church as the one, holy, catholic and apostolic Church, of which the
Nicene Creed speaks. From the Catholic point of view, the same
self-awareness applies: the one, holy, catholic and apostolic Church
``subsists in the Catholic Church'' (\emph{Lumen Gentium}, 8); this does
not exclude acknowledgement that elements of the true Church are present
outside the Catholic communion.
